% !PRINTED 161
\weberpage{161}

If therefore we set
\weq{3}{\dfrac{\sigma'(u)}{\sigma(u)} = \zeta(u)}
then $\zeta(u)$ is a single-valued function of $u$, and indeed
an elliptic integral of the second kind:
\weq{4}{\zeta(u) = -\int \wp(u)\,du = -\int \dfrac{\wp\,d\wp}{\sqrt{4\wp^3 - g_2\wp - g_3}} \ [\S\,46,\,(19)],}
in which the additive constant is so determined that $\zeta(u)$
must be an \gesperrt{odd} function. The periodicity of the $\zeta$-function
follows from the period equations of the $\sigma$-function
[\S\,35, (9)]:
\weq{5}{\begin{aligned} \zeta(u+\omega_1) - \zeta(u) &= 2\eta_1, \\ \zeta(u+\omega_2) - \zeta(u) &= 2\eta_2. \end{aligned}}
Thus $2\eta_1$ and $2\eta_2$ are represented as complete integrals of the second
kind (analogous to the $E$, $E'$):
\weq{6}{2\eta_1 = -\int\limits_u^{u+\omega_1} \wp(u)\,du, \quad 2\eta_2 = -\int\limits_u^{u+\omega_2} \wp(u)\,du}
% [sic] The German source duplicates the article in “mit der der Legendre[schen] Relation”; English conveys the intended syntax.
with the equation corresponding to the \gesperrt{Legendre} relation
\weq{7}{\eta_1\omega_2 - \eta_2\omega_1 = \pi i \ [\S\,35,\,(10)].}
In order to form the addition theorems for the \gesperrt{Weierstrass}
transcendents, we start from the $\sigma$-function.
For this function we obtain from (1), using the addition formula
for $\vartheta_{11}$, \S\,22, (4):
\weq{8}{\dfrac{\sigma(u+v)\sigma(u-v)}{\sigma(u)^2\sigma(v)^2} = \wp(v) - \wp(u),}
and from this, by logarithmic differentiation:
\weq{9}{\begin{aligned} \zeta(u+v) + \zeta(u-v) - 2\zeta(u) &= \dfrac{\wp'(u)}{\wp(u) - \wp(v)}, \\ \zeta(u+v) - \zeta(u-v) - 2\zeta(v) &= \dfrac{-\wp'(v)}{\wp(u) - \wp(v)}, \end{aligned}}
\weq{10}{\zeta(u+v) = \zeta(u) + \zeta(v) + \dfrac{1}{2}\dfrac{\wp'(u) - \wp'(v)}{\wp(u) - \wp(v)}.}
From this one derives, by a further differentiation, the
addition theorem for the $\wp$-function. From (9) one obtains,
by differentiation with respect to $u$ and $v$:
%CONT>>
